\SourceLeaf{321}%
or to another set (among the four variants that belong together in each case), one discovers only afterwards,\\\hspace*{2em}unless one has applied my criterion from the outset.

I do not give the criterion itself here, so as not to deprive you of the pleasure of finding it yourself. In any event, I hope I have made it sufficiently clear in the foregoing \textls[120]{what is at issue}. It is probably unnecessary to remark that the whole subject is highly elementary. That makes it all the more surprising that it has been so completely overlooked.

\mbox{. . . . . . . . . .}

\Needspace{12\baselineskip}\begin{center}[2.]\end{center}

\begin{center}\textsc{Gauss} to \textsc{Gerling}. Göttingen, 24 October 1840.\end{center}

\indent \mbox{. . . . . . . .} I read your little work on \textsc{Pothenot}'s problem with pleasure. It will be very instructive and useful to all who wish to apply the problem in practice. But I have often wondered that the many authors who have dealt with it have altogether passed over a part that is quite essential from a theoretical point of view, namely stating the mathematical conditions under which the problem is \textls[120]{physically possible} and presenting those conditions in a simple, elegant form. I think I can make myself understood most easily as follows.

Applied to given numerical data, the rules of calculation always yield one single determinate result (unless the four points lie on a single circle, in which case one obtains $\frac{0}{0}$). But one can readily see that one obtains \textls[120]{the same} result whether the second point is required to lie to the right of the first by the angle $A$ or to its left by $180^0-A$. The same holds for the third point. One can put this another way.

If it is prescribed that the second point lie to the right of the first by $A$, and the third to the right of the first by $B$, one finds a result \textls[120]{identical} to what one would find for three other sets of data, namely
